\documentclass[10pt]{article}
\usepackage[margin=1.02in]{geometry}
\usepackage{amsmath,amssymb,amsthm,mathtools}
\usepackage{graphicx}
\usepackage{xcolor}
\usepackage[hidelinks]{hyperref}
\usepackage{microtype}

\newtheorem{theorem}{Theorem}
\newtheorem{corollary}{Corollary}
\newcommand{\R}{\mathbb R}
\newcommand{\Lip}{\operatorname{Lip}}
\newcommand{\pa}{\mathrm{PA}_{\mathrm{fin}}}

\definecolor{statusred}{RGB}{151,37,32}

\title{An Oscillatory Shear Excludes the\\
Lipschitz Approximation Endpoint}
\author{fa\_banach\_001 / agent\_lane\_14}
\date{13 August 2026}

\begin{document}
\maketitle

\noindent
\fcolorbox{statusred}{statusred!6}{%
  \begin{minipage}{0.92\textwidth}
  \textbf{Status: candidate full counterexample, likely valid; human review
  needed.}  The endpoint $1$ does not belong to the approximation set
  $B(1,1)$ of Bellido--Mora-Corral.  In fact, one explicit bi-Lipschitz shear
  stays Lipschitz distance at least one from every map affine on a finite
  triangulation.  Together with the source theorem, this gives
  $B(1,1)=(0,1)$.
  \end{minipage}}

\section{Source question}

Bellido and Mora-Corral study approximation of planar H\"older
homeomorphisms by piecewise affine homeomorphisms over Euclidean finite
complexes \cite{BMC}.  For $0<\alpha,\widetilde\alpha\leq1$, they let
$B(\alpha,\widetilde\alpha)$ be the set of exponents $0<\beta\leq1$ for
which, on every closed polygon $\Omega$, piecewise affine homeomorphisms are
dense in
\[
 \bigl\{h\in C^\alpha(\Omega,\R^2):h\text{ is a homeomorphism and }
 h^{-1}\in C^{\widetilde\alpha}(h(\Omega),\R^2)\bigr\}
\]
in the $C^\beta$ norm.  Their Theorem~7.6 implies
$\sup B(1,1)=1$.  On PDF page~45 they ask whether the endpoint belongs:
\[
 1\stackrel{?}{\in}B(1,1).
\]

\begin{figure}[ht]
  \centering
  \includegraphics[width=\textwidth]{figures/open_problem_crop.png}
  \caption{The endpoint question on PDF page~45 of arXiv:0806.3366.}
\end{figure}

\section{Definitions}

For a map $F$ on the unit square $Q=[0,1]^2$, write
\[
 |F|_1=\sup_{p\neq q\in Q}\frac{|F(p)-F(q)|}{|p-q|},
 \qquad \|F\|_1=\|F\|_\infty+|F|_1.
\]
This is the source's exponent-one H\"older norm.  Let $\pa(Q,\R^2)$
denote all maps affine on the triangles of some finite triangulation of $Q$;
injectivity is not imposed in this larger class.

\section{Proof intuition}

The endpoint norm detects slopes rather than only function values.  A finite
piecewise affine map has only finitely many slopes along each boundary edge,
so near a corner its boundary slope is constant.  We build a horizontal
shear whose boundary slope continues to alternate between $-1$ and $+1$ at
arbitrarily small scales.  The shear and its inverse remain Lipschitz because
the slope is bounded.  No constant can be within less than one of both
$-1$ and $+1$, which yields a uniform gap from every finite piecewise affine
map.

\section{Counterexample and exact endpoint}

Define $g:[0,1]\to\R$ by
\[
 g(0)=0,\qquad g(t)=t^2\sin(1/t)\quad(t>0),
\]
and define
\begin{equation}
 h:Q\longrightarrow\R^2,\qquad h(x,y)=(x+g(y),y).
 \label{eq:shear}
\end{equation}

\begin{theorem}
\label{thm:gap}
The map $h$ in \eqref{eq:shear} is a bi-Lipschitz homeomorphism from $Q$
onto $h(Q)$.  Moreover, for every $f\in\pa(Q,\R^2)$,
\begin{equation}
 \|h-f\|_1\geq |h-f|_1\geq1.
 \label{eq:gap}
\end{equation}
Consequently, $1\notin B(1,1)$.
\end{theorem}

\begin{proof}
For $t>0$,
\begin{equation}
 g'(t)=2t\sin(1/t)-\cos(1/t),
 \label{eq:derivative}
\end{equation}
so $|g'(t)|\leq3$ on $(0,1]$.  Also $g'(0)=0$ directly from the
difference quotient.  Thus $g$ is Lipschitz.  The shear is injective because
it preserves the second coordinate, and its inverse on $h(Q)$ is
\[
 h^{-1}(u,v)=(u-g(v),v).
\]
The triangle inequality and the Lipschitz bound for $g$ show that both maps
are Lipschitz.  Hence $h$ satisfies the defining regularity hypotheses for
$B(1,1)$.

Fix $f\in\pa(Q,\R^2)$ and restrict it to the left edge
$E=\{(0,t):0\leq t\leq1\}$.  Intersecting a finite triangulation with $E$
gives a finite subdivision of $E$, so there are $\delta>0$ and $a,c\in\R$
such that
\[
 \pi_1 f(0,t)=a+ct\qquad(0\leq t\leq\delta).
\]
The first coordinate of $(h-f)(0,t)$ on this interval is therefore
$d(t)=g(t)-a-ct$.  For all sufficiently large $n$, both
\[
 t_n=\frac1{2\pi n},\qquad s_n=\frac1{(2n+1)\pi}
\]
belong to $(0,\delta)$.  Equation \eqref{eq:derivative} gives exactly
\[
 g'(t_n)=-1,\qquad g'(s_n)=1.
\]
The absolute value of a derivative at any differentiability point is bounded
by the Lipschitz seminorm.  Consequently,
\begin{align*}
 |h-f|_1
 &\geq \Lip(d;[0,\delta])\\
 &\geq \max\{|{-1-c}|,|{1-c}|\}
  =1+|c|\geq1,
\end{align*}
which proves \eqref{eq:gap}.  Since this excludes approximation even by the
larger class $\pa(Q,\R^2)$, it certainly excludes finite piecewise affine
homeomorphisms.  The square is a closed polygon, so the universal density
property defining $1\in B(1,1)$ fails.
\end{proof}

\begin{corollary}
The source approximation set is exactly
\[
 B(1,1)=(0,1).
\]
\end{corollary}

\begin{proof}
For $\alpha=\widetilde\alpha=1$, source Theorem~7.6 gives approximation in
every exponent $0<\beta<1$, hence $(0,1)\subseteq B(1,1)$.  Theorem
\ref{thm:gap} excludes the only remaining endpoint.
\end{proof}

\section{Verification, scope, and novelty}

The included script \texttt{code/verify\_shear.py} checks 3,500 instances of
the two exact slope sequences against seven affine slopes and verifies the
shear/inverse identities on a $101\times101$ grid.  This is only a numerical
sanity check; the proof is the elementary calculation above.  Details are in
\texttt{verification.md}.

The counterexample settles precisely the source's $B(1,1)$ endpoint.  It is
stronger than needed because the lower bound holds for noninjective finite
piecewise affine maps as well.  It makes no assertion about approximation by
countably piecewise affine maps with subdivisions accumulating at the corner,
nor about stronger quantitative bi-Lipschitz approximation questions.

Cheap run indexes, the parsed local corpus, and bounded web/arXiv searches
through 13 August 2026 used the exact endpoint phrase, title and authors,
$B(1,1)$, finite piecewise affine density, Lipschitz-norm approximation, and
oscillatory-shear terms.  They found the source and later Sobolev and
bi-Lipschitz approximation literature, but no explicit answer to this exact
endpoint question.  See \texttt{novelty\_search.md}.  This is not an
exhaustive priority claim.

\section{Human-review recommendation}

The highest-value check is the finite-complex convention: the restriction of
a map affine on a finite triangulation to a polygon edge must have a first
affine interval.  A reviewer should also confirm that the source's
$\|\cdot\|_1$ is the ordinary Lipschitz norm.  Once those definitions are
matched, the explicit inverse and the two derivative sequences make the
counterexample immediate.

\begin{thebibliography}{9}

\bibitem{BMC}
J. C. Bellido and C. Mora-Corral,
\emph{Approximation of H\"older continuous homeomorphisms by piecewise affine
homeomorphisms}, arXiv:0806.3366 (2008).

\end{thebibliography}

\end{document}
